\section{Study Comparison Following Webster and Watson}
Following the concept matrix approach proposed by \parencite{websterAnalyzingPrepareFuture2002}, this section presents a structured comparison of the 13 identified studies on wearable seizure detection and forecasting. The matrix organizes studies along key dimensions relevant to clinical translation and technical implementation.

\subsection{Concept Matrix: Study Overview}

Table~\ref{tab:study-matrix} provides a detailed comparison of all included studies, organized by primary objective, methodology, and outcomes. The matrix allows cross-study analysis of trends in sensor modalities, algorithmic approaches, and validation strategies across the 2013-2026 period.

\begin{landscape}
\begin{table}[htbp]
\centering
\small
\caption{Webster \& Watson Concept Matrix: Wearable Seizure Detection and Forecasting Studies (2013-2026)}
\label{tab:study-matrix}
\setlength{\tabcolsep}{3.5pt}
\renewcommand{\arraystretch}{1.15}
\resizebox{\linewidth}{!}{%
\begin{tabular}{@{} l l l l l l >{\centering}p{0.7cm} >{\centering}p{0.6cm} >{\centering}p{0.7cm} >{\raggedright\arraybackslash}p{2.5cm} l l l >{\raggedright\arraybackslash}p{2.8cm} @{}}
\toprule
Study & Design & Sample & Device & Loc. & \thead{Para.} & \thead{Mod.} & \thead{Pers.} & Algorithm & Sens & Spec & FPR & Key Limitation \\
\midrule

\textbf{\underline{Detection Studies}} \\
\addlinespace

\textbf{Spahr et al. 2025} & Prospective & n=384, 49 CSs (103 total) & Empatica E4 (ACC) & Wrist & Ens & 1 & Glb & Ensemble 1D CNN (30 models) & 96\% & -- & $<$0.0125/h & EMU, offline \\
\addlinespace

\textbf{Reintjes et al. 2025} & Retro & n=120, 856 szs & Single-lead ECG & Chest & Anom & 1 & Subj & Anomaly: Matrix Profile, MADRID, TimeVQVAE & 38.0\%/98.16\% & -- & 1.91--39.75/h & Severe sens/FAR trade-off \\
\addlinespace

\textbf{Fine 2025} & Phase 1 & n=15/3 & 6-axis band (ACC+Gyro) & Wrist & Feat & M & Glb & ANN (594 handcrafted feat) & 100\%Test, 96\% complete & -- & 0.023/h & Small test, Possible overfitting \\
\addlinespace

\textbf{Dong et al. 2026} & Prospective & n=68, 1846 szs & NightWatch & Arm & 2-stg & M & Glb & CNN-LSTM + Attention & 71.6\% & -- & 0.165/h & Single-center \\
\addlinespace

\textbf{Wang et al. 2025} & -- & n=28, 62 szs & Biovital-P1 & Wrist & Feat & M & Glb & LSTM (40 hidden) & 56.40\% to 95.3\% & -- & 0.354/h & Hospital only \\
\addlinespace

\textbf{Singh Rathore 2024} & Case study & 1 pt, 6 hrs, 3.2M pts & Wearable (EDA+ACC+HR) & -- & Feat & M & Glb & MLP (multi-modal features) & 96.8\% & 94.8\% & -- & Single patient \\
\addlinespace

\textbf{Elemam et al. 2025} & Cross-sect & Q:198, HRV:30 & Camera PPG & Thumbs & Hyb & M & Glb & CNN + Rule-based fusion & 95.1\% Q, 93-90\% HRV, 92.5\% AUD & 97.06\% Q & -- & Small samples \\
\addlinespace

\textbf{Borujeny 2013} & 3-pt & 3 pts, 20 szs & MICAz ACC & Arm/thigh & Feat & 1 & Glb & KNN k=5 (time-domain feat) & 100\% & -- & 0 & Very small \\
\addlinespace

\textbf{\underline{Forecasting Studies}} \\
\addlinespace

\textbf{Vieluf et al. 2025} & Retro & n=70, 5437 d & Embrace & Wrist & Feat & M & Mix & DNN + harmonic features + diary & 82\% & 67\% & -- & Daily res only \\
\addlinespace

\textbf{Meisel et al. 2020} & LOSO CV & n=69, 452 szs & Empatica E4 & Wrist/Ankle & E2E & M & Glb & LSTM + 1D Conv & 51.2\% & -- & TiW 43.7\% & Pediatric \\
\addlinespace

\textbf{Stirling 2021} & Retro+pseudo & n=11, 136 szs & Fitbit & Wrist & Feat & M & Pat & LSTM+RF+LR (HRV feat) & -- & -- & AUC 0.74 & Self-report \\
\addlinespace

\textbf{Nasseri 2021} & Retro & n=6 & Empatica E4 & Wrist & E2E & M & Tmp & LSTM 4-layer (128 hidden) & AUC 0.80 (SD 0.15) & -- & TiW 0.9--7.2h/d & RNS required \\
\addlinespace

\textbf{Ode et al. 2023} & Retro & n=66, 85 szs & ECG & -- & Anom & 1 & Pat & Self-Attentive AE (attention) & 74\% & -- & 0.85/h & FPs in some pts \\
\bottomrule
\end{tabular}%
}
\par\medskip
\footnotesize
\textbf{Note:} ACC = accelerometer, ANN = artificial neural network, AE = autoencoder, CNN = convolutional neural network, CS = convulsive seizure, DNN = deep neural network, ECG = electrocardiogram, EDA = electrodermal activity, HR = heart rate, HRV = HR variability, LOSO = leave-one-subject-out, LSTM = long short-term memory, PPG = photoplethysmography, RRI = R-R interval, Sens = sensitivity, Spec = specificity, sz(s) = seizure(s), FPR = false positive rate, TiW = time in warning, -- = not reported. \textbf{New columns:} Para. = learning paradigm (Feat = feature-based with handcrafted features, E2E = end-to-end learning from raw data, Hyb = hybrid, 2-stg = two-stage, Anom = anomaly detection, Ens = ensemble); Mod. = modality count (1 = single-modal, M = multi-modal); Pers. = personalization (Glb = global model, Pat = patient-specific, Mix = mixed, Tmp = temporal split, Subj = subject split).
\end{table}
\end{landscape}

\subsection{Architecture and Modality Deep-Dive}

Table~\ref{tab:detailed-metrics} provides a detailed technical comparison of deep learning architectures, biosignal modalities, and deployment factors across all studies. This table directly addresses the research question of how different architectures and modalities compare in achieving optimal sensitivity-FAR trade-offs for ambulatory seizure monitoring.

\begin{landscape}
\begin{table}[htbp]
\centering
\small
\caption{Architecture and Modality Deep-Dive: Technical Dimensions Across Studies}
\label{tab:detailed-metrics}
\setlength{\tabcolsep}{3.5pt}
\renewcommand{\arraystretch}{1.1}
\resizebox{\linewidth}{!}{%
\begin{tabular}{@{} l >{\raggedright}p{1.8cm} >{\raggedright}p{2.2cm} >{\raggedright}p{2.0cm} >{\centering}p{1.2cm} >{\raggedright}p{1.8cm} >{\raggedright}p{2.0cm} >{\raggedright}p{1.8cm} @{}}
\toprule
\textbf{Study} & \textbf{Paradigm} & \textbf{Arch. Details} & \textbf{Modality Details} & \textbf{Fusion} & \textbf{Personalization} & \textbf{Trade-off Config} & \textbf{Deployment} \\
\midrule

\textbf{Detection Studies} \\
\addlinespace

\textbf{Spahr et al. 2025} & Ensemble & 30 1D CNN models, 14 conv layers, quantile & ACC 32~Hz, Euclidean norm, 30~s & -- & Global, tunable & 86--100\% @ $<$0.0125/h & Real-time (112~ms), on-device \\
\textbf{Reintjes et al. 2025} & Anomaly & Matrix Profile, MADRID, TimeVQVAE-AD & ECG 256$\to$8~Hz, bandpass & -- & Subject split & 38.0--98.2\% @ 1.91--39.75/h & Offline, hospital \\
\textbf{Fine 2025} & Feature-based & ANN, 594 handcrafted features & ACC+Gyro 6-axis, 10~s intervals & Early & Global, hold-out & 100\% @ 0.023/h & Offline PC, EMU \\
\textbf{Dong et al. 2026} & Two-stage & Pre-screening + CNN-LSTM + Attention & ACM 11-12$\to$20~Hz, PPG 100$\to$20~Hz, 5~min & Early & Global, 10-fold CV & 71.6\% @ 0.165/h & Real-time, NightWatch, home \\
\textbf{Wang et al. 2025} & Feature-based & LSTM (40 hidden) + ReLU + FC & ACC/GYR 50~Hz, SEMG 200~Hz, EDA 4~Hz, 4~s & Early & Global, hold-out & 56.4--95.3\% @ 0.354/h & Real-time, hospital \\
\textbf{Singh 2024} & Feature-based & MLP, multi-modal features & EDA, ACC, HR/HRV, 25k points & Early & Global, single pt & 96.8\% @ 94.8\% prec & Real-time, cloud? \\
\textbf{Elemam 2025} & CNN-based & CNN for HRV, CNN for Audio (separate) & PPG 250~Hz, Audio 10~s & No fusion & Global, unclear & 95.1\% @ 97.1\% spec & Real-time, hospital \\
\textbf{Borujeny 2013} & Feature-based & KNN k=5, time-domain features & ACC 2D 3~Hz, 9~s & -- & Global, 3 pts & 90--100\% @ 0--15\% & Real-time, server (316~mW) \\
\addlinespace

\textbf{Forecasting Studies} \\
\addlinespace

\textbf{Vieluf et al. 2025} & Hybrid & DNN + harmonic modeling (24-h) & EDA 4~Hz, ACC 32~Hz, Temp 1~Hz, 10~min & Early & Mixed, 5-fold + LOO & 82\% sens, 67\% spec @ F1 0.81 & Offline MATLAB, home \\
\textbf{Meisel et al. 2020} & End-to-end & LSTM only (10 units) & EDA/ACC/BVP/Temp $\to$4~Hz, 30~s & Early & Global, LOSO & 51.2\% @ IoC 14.1\% & Offline, hospital \\
\textbf{Stirling 2021} & Feature-based & LSTM+RF+LR ensemble, LR combines & HR 5~s, steps/sleep 1~min, hourly/daily & Decision & Patient-specific, weekly & AUC 0.74, 100\% pts & Home (Fitbit), app \\
\textbf{Nasseri 2021} & End-to-end & LSTM 4-layer, 128 hidden + FFT channels & ACC/BVP/EDA/Temp/HR $\to$128~Hz, 60~s & Early (17-ch) & Temporal split & AUC 0.80 (SD 0.15) @ TiW 0.9--7.2h/d & Offline, cloud, home \\
\textbf{Ode 2023} & Anomaly & Self-Attentive AE (attention) & ECG (RRI only), 45~s & -- & Patient-specific (99\% CL) & 74\% @ 0.85/h (99\% CL) & Real-time, cloud, hospital \\
\bottomrule
\end{tabular}%
}
\par\medskip
\footnotesize
\textbf{Note:} Paradigm: Feature-based = handcrafted features, End-to-end = learns from raw data, Hybrid = combines both, Ensemble = multiple models, Two-stage = sequential processing, Anomaly = unsupervised detection, CNN-based = convolutional approach. Arch. Details: architecture-specific components. Fusion: -- = single-modal, Early = feature-level, Decision = classifier output, No fusion = parallel separate models. Personalization: Global = population model, Patient-specific = individualized, Temporal = time-split within patient, LOSO = leave-one-subject-out, LOO = leave-one-out, Subject split = train/test separated by subject. Trade-off Config: operating points with sensitivity @ FAR. Deployment: real-time capability, processing location, setting. ACC = accelerometer, ECG = electrocardiogram, PPG = photoplethysmography, EDA = electrodermal activity, BVP = blood volume pulse, HR = heart rate, HRV = HR variability, SEMG = surface electromyography, TEMP = temperature, FFT = fast Fourier transform, CL = confidence limit, AUC = area under curve, TiW = time in warning, IoC = Improvement over Chance, EMU = epilepsy monitoring unit.
\end{table}
\end{landscape}

\subsection{Dimensional Analysis}

\subsubsection{Sensor Modality Trends}
\begin{itemize}
  \item \textbf{Wrist-worn accelerometry} dominates (8/13 studies), leveraging commercial devices (Empatica E4, Fitbit, Embrace).
  \item \textbf{Cardiac signals} (ECG, PPG, HRV) show increasing adoption for preictal detection, with 6 studies incorporating autonomic measures.
  \item \textbf{Multi-modal fusion} (ACC + EDA + PPG) appears in 4 recent studies, reflecting recognition that single modalities have limited specificity.
\end{itemize}

\subsubsection{Algorithmic Evolution}
\begin{itemize}
  \item \textbf{2013:} Traditional ML (KNN, ANN) with handcrafted features \parencite{borujenyDetectionEpilepticSeizure2013}.
  \item \textbf{2020-2021:} Emergence of LSTM and ensemble methods \parencite{meiselMachineLearningWristband2020,stirlingForecastingSeizureLikelihood2021}.
  \item \textbf{2023-2026:} Deep learning dominance with CNNs, attention mechanisms, and anomaly detection approaches \parencite{reintjesECGBasedDetectionEpileptic2025,odeDevelopmentEpilepticSeizure2023,spahrDeepLearningbasedDetection2025}.
\end{itemize}

\subsubsection{Validation Rigor}
\begin{itemize}
  \item \textbf{Patient-level preservation:} Only 4 studies use LOSO or patient-independent splits.
  \item \textbf{Prospective validation:} Only 3 studies \parencite{spahrDeepLearningbasedDetection2025,dongTwoLayerEnsembleMethod2022,elemamAutomatedValidatedTool2025} include prospective phases.
  \item \textbf{Real-world testing:} Limited to 2 studies \parencite{vielufSeizureMonitoringCombined2025,nasseriAmbulatorySeizureForecasting2021}, with most remaining in EMU settings.
\end{itemize}

\subsubsection{Performance Metrics Reporting}
Table~\ref{tab:metrics-summary} summarizes the heterogeneity in performance metric reporting, complicating cross-study comparison. Sensitivity ranges from 51.2\% (forecasting mean) to 100\% (tonic detection), while FPR reporting varies from per-hour to per-night, with several studies omitting this critical metric entirely.

Table~\ref{tab:detailed-metrics} reveals additional clinically relevant metrics that are critical for real-world deployment but inconsistently reported:
\begin{itemize}
  \item \textbf{Detection latency} is reported by only 4/13 studies, despite being critical for timely intervention.
  \item \textbf{Patient-level success rate} (proportion achieving usable performance) is reported by only 6/13 studies, ranging from 43.5\% to 100\%.
  \item \textbf{Precision/PPV} is rarely reported (3/13), limiting understanding of positive predictive value for clinical decision-making.
  \item \textbf{Generalization metrics} (LOSO validation, patient-independent testing) show that only Meisel et al. used true LOSO CV, with most studies reporting aggregated performance that may overestimate real-world utility.
\end{itemize}

\begin{table}[htbp]
\centering
\small
\caption{Performance Metrics Reporting Across Studies}
\label{tab:metrics-summary}
\resizebox{\textwidth}{!}{
\begin{tabular}{lccc}
\toprule
\textbf{Metric} & \textbf{Detection (n=9)} & \textbf{Forecasting (n=4)} & \textbf{Reported By} \\
\midrule
Sensitivity/Recall & 8/9 & 2/4 & Reintjes, Fine, Spahr, Dong, Wang, \\
& & & Singh Rathore, Elemam, Borujeny, Meisel, Ode \\
Specificity & 1/9 & 0/4 & Elemam only \\
FPR/FA rate & 7/9 & 2/4 & Spahr, Reintjes, Fine, Dong, Wang, \\
& & & Ode, Meisel, Nasseri, Stirling \\
AUC-ROC & 1/9 & 2/4 & Reintjes, Nasseri, Stirling \\
Detection Latency & 3/9 & 1/4 & Spahr, Fine, Nasseri \\
Patient Success Rate & 2/9 & 4/4 & Meisel, Nasseri, Ode, Stirling, Vieluf, Wang \\
Precision/PPV & 2/9 & 1/4 & Ode, Singh Rathore, Elemam, Wang \\
IoC / Time in Warning & 0/9 & 2/4 & Meisel, Nasseri \\
\bottomrule
\end{tabular}
}
\end{table}

\subsection{Synthesis and Gaps}

The concept matrix and detailed metrics reveal several patterns:

\begin{enumerate}
  \item \textbf{Clinical translation gap:} Despite high sensitivities ($>$90\% in several detection studies), FPR remains clinically unacceptable for standalone deployment. Only Dong et al. (0.165/h) and Fine (0.023/h) approach the ILAE Phase 3 benchmark for validated devices (see Section~\ref{sec:ilae-standards}).

  \item \textbf{Personalization vs. generalization:} LOSO approaches achieve lower but more realistic performance (51.2\% sensitivity, 43.5\% patient success) compared to patient-specific tuning. Meisel et al.'s LOSO CV results suggest that aggregated performance metrics may overestimate real-world utility.

  \item \textbf{Patient-level heterogeneity:} Patient success rates vary dramatically (43.5-100\%), with Meisel and Ode showing that fewer than half of patients achieve usable performance. Nasseri (5/6) and Stirling (100\%) demonstrate that responder identification remains a critical challenge.

  \item \textbf{Diary integration:} \parencite{vielufSeizureMonitoringCombined2025} demonstrates that hybrid approaches (wearable + clinical data) may outperform wearables alone, suggesting a direction for ambulatory systems where wearable-only performance is insufficient.

  \item \textbf{Metric heterogeneity:} Inconsistent reporting (sensitivity vs. accuracy vs. AUC vs. IoC) impedes meta-analysis and clinical benchmarking. Key clinical metrics like detection latency and patient success rate are reported by fewer than half of studies.
\end{enumerate}
